\scanpage{309}
\mspage{259}

\uline{Définition} Quand ces conditions sont satisfaites, on dit que le carré $Q$ est $W$-cartésien.

\uline{Corollaire 1} \deleted{Supposons que $W$ soit \unsure{forte-}} \inserted{Si $Q$ est $W$-cartésien, alors} \inserted{l'image de} $Q$ \inserted{dans $\mathrm{Ho}_W(\Psi^{\sim})$} est dans l'image essentielle de \deleted{\unsure{$\varepsilon$|$i$}}
\[
\unsure{j|i}_*\colon\mathrm{Ho}_W(\Psi)\longrightarrow\mathrm{Ho}_W(\Psi^{\sim}),
\]
i.e. la flèche canonique dans $\mathrm{Ho}_W(\Psi^{\sim})$
\[
\deleted{\unsure{i_*|i^*}}\,Q\longrightarrow i_*i^*Q
\]
est iso. La réciproque est vraie si $W$ est fortement \unsure{saturée|saturé}.

\uline{Corollaire 2} Considérons un diagramme comm. dans $M$
\diagram{\(\text{(D)}\quad\begin{tikzcd}[ampersand replacement=\&,row sep=small,column sep=small]
X'_0\arrow[dd]\&\&X'_1\arrow[ll]\arrow[dd]\&\&X'_2\arrow[ll]\arrow[dd]\\
\&1\&\&2\&\\
X_0\&\&X_1\arrow[ll]\&\&X_2\arrow[ll]
\end{tikzcd},\)}
d'où deux carrés commutatifs 1, 2. et un carré composé \unsure{3},
\diagram{\begin{tikzcd}[ampersand replacement=\&,row sep=small,column sep=small]
X'_0\arrow[dd]\&\&X'_2\arrow[ll]\arrow[dd]\\
\&3\&\\
X_0\&\&X_2\arrow[ll]
\end{tikzcd}.}

Supposons le carré 1 $W$-cartésien. Alors le carré 2 est $W$-cartésien ssi le carré 3 l'est.

\uline{Dém} Considérons une factorisation de
\diagram{\begin{tikzcd}[ampersand replacement=\&,row sep=small]
X'_0\arrow[d]\\X_0
\end{tikzcd}}
en
\diagram{\begin{tikzcd}[ampersand replacement=\&,row sep=small,column sep=small]
\overline X'_0\arrow[dr]\&X'_0\arrow[l]\arrow[d]\\
\&X_0
\end{tikzcd}}
avec $(X'_0\mathrel{\unsure{-|\to}}\overline X'_0)\in W$, et $(X'_0\mathrel{\unsure{-|\to}}X_0)\in\mathrm{Fib}_W$. On peut \unsure{donc|alors} remplacer dans (D) $X'_0$ par $\overline X'_0$, \deleted{\unsure{en s}} obtenant un diagramme $\overline D$, et on voit \unsure{aussitôt}
% archive sheet number: 309
